\section{Holy Rome, Eternal Rome}

Before translating Evola's review of a book by Maurras, we plan to provide a chapter from \textbf{Guido de Giorgio}'s La Tradizione Romana (The Roman Tradition) called ``Dante and the Holy Culmination of the Roman Tradition.''

De Giorgio is most certainly the only man to have collaborated with both \textbf{Rene Guenon} and \textbf{Julius Evola}. In his early life, he lived in Tunisia for a time where he was initiated into Islamic esoterism. He then relocated to Paris where he worked with Guenon on two Traditional periodicals. Back in Rome, he participated in the Ur magic group, along with Evola, Duke Colonna de Cesaro, and several others. He also worked with Evola on the political journal La Torre. Suddenly, he abandoned public life, retiring to the Italian Alps where he lived an ascetic and contemplative life. His books were never published in his lifetime.

What is of interest here is his understanding of the European Middle Ages, since it is the traditional society that is closest to us in time and culture. However, a precise understanding is a matter of contention among the main writers of Tradition. If Guenon is the thesis, then Evola is the antithesis and de Giorgio the synthesis.

\paragraph{Rene Guenon}


Guenon was not interested in the ancient Western Tradition, since for him it is defunct. He does, however, recognize a traditional metaphysics in both neo-Platonism and Aristotelianism. He did consider the religious and social systems of the Middle Ages to be Traditional. He points to the metaphysical teachings of the time and he infers the existence of legitimate initiatic orders based on certain writings, especially in Dante. On this basis, we can assuredly assume there was an intellectual elite behind that system, since Tradition does not come about by chance, nor from below. In particular, it was not the creation of slaves and an underclass.

\paragraph{Julius Evola}

For Evola, the Christian Middle Ages arose out of the wreckage of ancient Rome, already in decline in what he called Semitic Rome. It was a lower class movement, although it was at some point taken over and redirected by a new elite, which was composed of the intellectual meeting of Roman and Nordic elements. This elite managed to restore traditional elements to a large extent, though not completely. Unfortunately, this project was made difficult due to recurrent conflicts between the religious and social spheres, or more specifically, between the Church and Empire.

In the dozen or so years between the publication of Pagan Imperialism\footnote{\url{http://www.gornahoor.net/PaganImperialism/PaganImperialism.htm}} and On the Subterranean History of Rome\footnote{\url{http://www.gornahoor.net/?p=3832}}, Evola softened his attitude toward the Medieval church to some extent. For him, it served as a bulwark against contemporary neo-pagan and neo-spiritualist movements, with their pseudo-traditional influence. Against Cesaro, he wrote:

\begin{quotex}
we see Christianity becoming Roman with Catholicism: purifying itself of its original anarchic, universalistic, and humanitarian aspects, and giving rise, in the Middle Ages, to a civilization that is characteristic of the type we articulated: hierarchical, tied to traditions of caste and blood, interspersed with initiatic elements
\end{quotex}


Specifically, he opposed the evolutionistic thrust behind neo-spiritual movements, which saw history as the progression from a primitive state to ever higher states culminating in the age of love. This he regarded as a heresy, both in its narrow religious sense, but also in the wider sense as a falsification and denial of traditional elements. So when \textbf{Alain de Benoist}, for example, considers modern liberalism to be the secularization of theology, he can only be referring to the heretical Christianized neo-spiritualism like that of Cesaro. It is absolute nonsense in terms of the traditional elements still alive in the Middle Ages.

\paragraph{Guido de Giorgio}


Whereas Guenon is uninterested in any historical relations between various traditions, de Giorgio instead sees a direct relationship between the \textbf{Ancient Rome} of the pagans and the \textbf{New Rome} of the Middle Ages. Taking Guenon quite seriously, he relies on \textbf{Dante} as an inspired poet and the Divine Comedy as a sacred poem. \textbf{Vergil} takes him from the spirituality of Ancient Rome and Beatrice to the New Rome.

As Evola pointed out, the legend of \textbf{Aeneas}, as told in Vergil's poem, represents the end of one world and the beginning of the new with the founding of Rome. We should understand Beatrice in a similar way: she shows Dante the New Rome after the destruction of the previous era. Furthermore, we want to understand this as the way to the \textbf{Supreme Identity} via the alchemical marriage.

De Giorgio's perspective almost regards pagan Rome as the true Old Testament, or at least a second Old Testament, to the New Testament of the New Rome. His insights come from a deep spiritual understanding of Rome. I hope this unusual piece will bring as much interest and commentary as our lighter pieces.


\flrightit{Posted on 2012-03-20 by Cologero}

\begin{center}* * *\end{center}

\begin{footnotesize}
\begin{sffamily}
\texttt{Boreas on 2012-03-21 at 09:07 said:}

De Giorgio seems to be very interesting as an author and in person. I'd really like to get my hands on his books some time, right now I have nothing but the english translations of the book `Introduction to Magic' in which he probably contributed in some way.

I think Guénon is un-surpassed in the field of esotericism and traditional metaphysical knowledge, but his personal disregard concerning the western esoteric tradition has recently caused him to fall below other more closer to myself personally. Just the other day we talked about the Mithraic mysteries, initiation and esotericism with a friend of mine, and it bumped to my mind that it was probably one of the only issues Guénon did not directly deal in his writings (correct me if I'm wrong); maybe because of this very same reason.

``De Giorgio's perspective almost regards pagan Rome as the true Old Testament, or at least a second Old Testament, to the New Testament of the New Rome.''

If perceived ``aionically'', I think all traditions before the coming of Christ can be seen as part of ``The Old Testament'' \& being under the impulse of Justice more than Love, with all their correspondences in culture, law and spiritual methodology. ``The Old Covenant'' as an esoteric concept is far wider than the 2000-3000 years of the Hebraic / Judaic Tradition, and the Old Testament as a book is a conglemarate of writings from several different traditions (Sumerian, Babylonian, Persian, Egyptian etc.), and in the hebraic versions some ancient myths curiously transform and turn upside down, maybe because of the guilt complex so prevalent in the judaic religious psyche, for example, when `the heroic act' transforms into `an unforgivable sin' (as described by Evola in the beginning of `The Hermetic Tradition' and elsewhere also).

I think the last mentioned fact is also the very key to the two thousand years of persecutions of alchemists and other torch-bearers of the Regal Tradition in the underground veins of the Christendom.

\hfill

\texttt{Cologero on 2012-03-21 at 12:14 said:}

De Giorgio's pseudonym was ``Havismat''; I don't think any of his writings were included in the English version of Introduction to Magic. I'll check the Italian version to see if there is anything of interest.

By the way, I'm trying to locate a copy of his ``Dio e il Poeta'', if any readers can help me out.

\hfill

\end{sffamily}
\end{footnotesize}

